\section{Related Work}

\subsection{Geospatial Foundation Models}
Recent years have seen rapid growth in foundation models for Earth observation. Prithvi-100M~\cite{jakubik2023prithvi}, SatMAE~\cite{cong2022satmae}, Scale-MAE~\cite{reed2023scalemae}, and SpectralGPT~\cite{hong2024spectralgpt} all extend the masked autoencoder paradigm~\cite{he2022mae} from natural images to multi-spectral and multi-temporal satellite imagery. These models share the Vision Transformer backbone~\cite{dosovitskiy2021vit} but differ from conventional vision encoders in three important ways: they ingest multi-spectral inputs rather than RGB, they exploit large unlabeled satellite corpora through self-supervised pre-training~\cite{manas2021seco,wang2022ssl4eo}, and they must support broad variation in spatial resolution and sensor configuration. While downstream transfer results are encouraging, the question of \emph{how} to adapt such models efficiently to new sensor combinations has received much less attention than the pre-training itself.

\subsection{Parameter-Efficient Fine-Tuning}
Parameter-efficient fine-tuning (PEFT) has become the standard recipe for adapting large pre-trained models without updating all parameters. BitFit~\cite{zaken2022bitfit} restricts learning to bias terms; LoRA~\cite{hu2022lora} injects low-rank updates into linear projections; Houlsby adapters~\cite{houlsby2019adapter} insert bottleneck modules between transformer blocks; and prompt-based methods~\cite{lester2021prompt,jia2022vpt} prepend learnable tokens to the input. Subsequent work has unified these strategies under a common framework~\cite{he2022unified} and explored composition across tasks~\cite{pfeiffer2020adapterfusion}. In language models and conventional vision transformers, these methods often approach full fine-tuning while reducing memory and storage costs by one to three orders of magnitude. However, their relative ranking depends strongly on architecture, task, and pre-training objective, which raises the possibility that PEFT rankings established in mainstream domains may not transfer to geospatial models.

\subsection{PEFT for Remote Sensing}
Although remote sensing research increasingly relies on pre-trained transformers, systematic evaluation of PEFT methods in this domain remains limited. Existing work tends to focus on a single adaptation method, a single dataset, or an application-specific pipeline; PEFT is typically adopted pragmatically rather than studied comparatively. Closely related lines of work, such as training-free adaptation in vision-language models~\cite{zhang2022tip}, show that adaptation strategy can dominate over backbone choice, but this observation has not been verified for geospatial foundation models. As a result, the field still lacks a controlled benchmark answering basic but practically important questions: Does LoRA~\cite{hu2022lora} work as expected on geospatial transformers? Which adapter family is most robust under heterogeneous input configurations? How much of the gap to full fine-tuning can be closed under small parameter budgets? Our work is designed to answer these questions directly.
